\scanpage{9}
\deleted{voisinage ouvert $W\ni t$, induisant ce qui est donné sur $W\cap U$. En fait, \deleted{d'après le} \inserted{si on dispose du} th. de Tate \deleted{\unsure{ceci}} \unsure{pas seulmt|pas seult} en \unsure{inégales|inégale} caractéristique, on pourrait prouver}

b) Cas $S$ \deleted{normal} \inserted{géom. unib.}, $A_U$ muni d'une rigidification de Jacobi. On a donc un morphisme
\margin{\unsure{R}}
\[
U\longrightarrow R
\]
\margin{schéma modulaire à échelon [utilisant aussi une polarisation de $A_U$]}
d'où par \inserted{l'adhérence du} le graphe un sous-schéma fermé $S'$ de $S\times R$. Le critère valuatif et le résultat de Raynaud dans le cas d'un trait impliquent que $S'\longrightarrow S$ est propre. Il reste à prouver que c'est un isom [tout revient en effet, grâce à Tate, à prouver que $A_U$ se prolonge en $A$, quand $T_p(A_U)$ se prolonge en un $M$ -- cet $M$ étant en effet unique à iso unique près). Par le Main Theorem, on est ramené à prouver que \deleted{les pts de $S'$} \inserted{ttes fibres} $S'_s$ \unsure{sont} \uline{discrètes}. Introduisons la normalisée $S''$ de $S'$, qu'on supposera d'abord \uline{finie} sur $S'$. Alors le Th. de Tate implique que le gpe $\unsure{b|p}$-divisible $T_p(\underbrace{A_{S''}}_{\inserted{\text{image inverse de }A_R}})$ est l'image inverse de $\deleted{\unsure{T_p\text{ de }A|T_p(A)}}M$. Donc la restriction de $A_{S''}|S''_s$, correspondant à un $T_p$ ``constant'', et $S''_s$ étant géom connexe,
\diagram{\begin{tikzcd}[ampersand replacement=\&,row sep=small,column sep=small]
 \& S''\arrow[d] \& U''\arrow[l,hook]\arrow[d,phantom,"\rotatebox{90}{$\simeq$}" description]\\
 R \& S'\arrow[l]\arrow[d] \& U'\arrow[l,hook]\arrow[d,phantom,"\rotatebox{90}{$\simeq$}" description]\\
 \& S \& U\arrow[l,hook]
\end{tikzcd}}

% archive sheet number: 8
